\documentclass[11pt]{book}

\usepackage{ppbook}

\begin{document}

\setcounter{chapter}{3}
\chapter{LLMs for Invariant Synthesis and Verification}

\section{LLM + invariants}

Main question = can LLMs generate useful program
invariants?

One approach used a synthetic Java dataset.

Daikon was used to get candidate invariants first.

\begin{ppdefn}
Daikon is a dynamic invariant generator.

Basically, it runs programs on test inputs and watches
the values of variables during execution.

Then it gives candidate properties which appeared to
hold over those executions.
\end{ppdefn}

The basic flow is

\[
\text{program + test runs}
\rightarrow
\text{variable values}
\rightarrow
\text{candidate invariants}
\]

These candidates were then used as training data
for the LLM.

The invariants are not only for loops.

Some types are:

\begin{itemize}
    \item object invariants
    \item class invariants
    \item function entry invariants
    \item function exit invariants
    \item loop invariants
\end{itemize}

Function entry $\rightarrow$ precondition.

Function exit $\rightarrow$ postcondition.

The work mainly looks at invariants around function
entry/exit and loop headers.

\section{Precision and recall}

For generated invariants,

\[
Precision = \frac{TP}{TP+FP}
\]

\[
Recall = \frac{TP}{TP+FN}
\]

Precision tells us how many predicted invariants
are actually correct.

Recall = how many of the relevant invariants were
actually found.

The fine-tuned model had about

\[
86\% \text{ precision and } 86\% \text{ recall}
\]

\begin{ppkey}
The important part is validation.

The LLM is only generating candidates.
The predicted invariant is checked using static analysis.
\end{ppkey}

So the flow is

\[
\text{LLM}
\rightarrow
\text{candidate invariant}
\rightarrow
\text{static analysis}
\]

not just directly trusting the generated statement.

Another observation:

Daikon gets worse when fewer execution traces
are available.

But a larger LLM with zero-shot prompting could
still find some invariants that Daikon had not found.

So LLM may discover new invariants rather than
just copying the ones produced by dynamic analysis.

\section{LEMUR}

Another work = LEMUR for automated program verification.

Basic idea is to use an LLM as an oracle for proposing
properties, while a verifier checks them.

Need to keep the distinction:

\[
\text{LLM proposes}
\]

\[
\downarrow
\]

\[
\text{verifier checks}
\]

\[
\downarrow
\]

\[
\text{proof succeeds / fails}
\]

There is an assumption associated with a program point.

If the assumption holds at that point, execution can
continue without changing the program.

If it does not hold, execution terminates.

A property that stays valid at the required program
location can act as an invariant.

\section{Proof process}

Configuration contains roughly

\[
\langle P,A,M\rangle
\]

where

\[
P = \text{program}
\]

\[
A = \text{assumptions}
\]

\[
M = \text{proof goals}
\]

Start with the program and the initial proof goal.

Then apply rules to change the configuration until
we get

\[
Success
\]

or

\[
Failure
\]

\section{Propose}

The oracle looks at the program + current proof goal
and suggests new properties.

The useful properties should:

\begin{itemize}
    \item be invariants
    \item help imply the proof goal
\end{itemize}

Initially these proposed properties are treated
as assumptions.

They still have to be proved to actually be invariants.

\section{Repair}

Sometimes the proposed assumption is not strong enough
or something goes wrong during verification.

Then another property can be generated.

The idea is

\[
\text{old assumption}
\rightarrow
\text{strengthen / correct it}
\]

The repair operation uses the current proof goal
and the previously proposed assumption.

If a property already appears in the proof trail
but the verifier finds it false, repair can also
be used to replace it with another property.

So it is not just propose once and stop.

\section{Decide}

Suppose the verifier proves

\[
q\Rightarrow p
\]

Then $q$ can become the next proof goal.

The system can use $k$-induction for this kind of
reasoning.

Basic idea of induction:

\[
\text{base case}
+
\text{inductive step}
\]

to establish that the property keeps holding.

\section{Backtrack}

Sometimes more than one rule can be applied.

Then the system may go back to the previous proof goal
and choose a different assumption.

So proof search is not always one straight path.

Roughly,

\[
\text{propose}
\rightarrow
\text{verify}
\rightarrow
\begin{cases}
\text{decide}\\
\text{repair}\\
\text{backtrack}
\end{cases}
\]

\section{Success / fail}

If the verifier establishes the required property,
we reach SUCCESS.

If the verifier establishes that the property does
not hold, then FAIL.

This gives an important soundness idea:

\[
SUCCESS \Rightarrow \text{property is an invariant}
\]

and

\[
FAIL \Rightarrow \text{property is not an invariant}
\]

provided the rule applications used are valid.

\section{Main point}

LLM is useful mainly for the search part.

It can suggest properties which may be hard for
the verifier to come up with itself.

But actual correctness comes from the verifier.

\begin{ppkey}
The basic setup is

\[
\boxed{
\text{LLM proposes}
\quad\rightarrow\quad
\text{verifier checks}
}
\]

LLM $\neq$ proof.
\end{ppkey}

For invariant synthesis,

\[
\text{program}
\rightarrow
\text{candidate invariants}
\rightarrow
\text{validation}
\]

For LEMUR,

\[
\text{proof goal}
\rightarrow
\text{propose property}
\rightarrow
\text{check}
\rightarrow
\text{repair / decide / backtrack}
\]

The useful thing is combining the LLM's ability to
generate possible properties with the formal verifier's
ability to check them.

\end{document}
